% !TEX TS-program = lualatex
% !TEX encoding = UTF-8 Unicode
\documentclass[ngerman,11pt]{article}
\usepackage{fontspec}
\usepackage[ngerman]{babel}
\usepackage[ngerman]{datetime2}
\usepackage[a4paper, total={16cm, 27cm}]{geometry}
\usepackage{amsmath}
\usepackage{amsfonts}
\usepackage{amssymb}
\usepackage{tcolorbox}
\usepackage{enumitem}
\usepackage{url}
\usepackage[colorlinks=true, linkcolor=blue, urlcolor=blue]{hyperref}
\usepackage{listings}
\usepackage{tikz}

\lstset{
    language=Python,
    basicstyle=\ttfamily\small,
    keywordstyle=\color{black},
    stringstyle=\color{black},
    commentstyle=\color{black},
    breaklines=true,
    frame=single,
    framesep=1mm,
    framextopmargin=1mm,
    framexbottommargin=1mm,
    showspaces=false,
    showstringspaces=false,
    numbers=left,
    numberstyle=\tiny\color{gray}
}
\date{\DTMgermanmonthname{\month} \the\year}

\begin{document}
\hrule
\begin{center}
{\large\bf Algorithmen und Funktionen}\\[-0mm]
Studienkolleg der TU Berlin\hfill Till Zoppke\\
Oberkurs Informatik \hfill \DTMgermanmonthname{\month} \the\year\\[2mm]
\end{center}
\hrule

%\par\bigskip

% ============================================================

\section*{Orangen Wiegen}
 Sie haben \textbf{neun Orangen}. Alle sehen gleich aus, aber \textbf{eine} ist schwerer als die anderen. 

\begin{enumerate}
\item Wiegen Sie die Orangen mit einer virtuellen Balkenwaage. Öffnen Sie dazu die Simulation im ISIS-Kurs.\footnote{\url{https://isis.tu-berlin.de/mod/resource/view.php?id=2253481}} Probieren Sie verschiedene Strategien aus.

\item Schreiben Sie eine Anleitung, wie man mit möglichst wenig Wiegevorgängen die schwere Orange finden kann. Die Anleitung soll eindeutig sein, sodass eine andere Person sie ohne Nachfragen ausführen kann.

\begin{tcolorbox}[colback=white, colframe=black, boxrule=0.3pt, height=7cm]
\vspace*{\fill}
\end{tcolorbox}

\item Finden Sie die Person mit dem gleichen Buchstaben in der linken oberen Ecke. Probieren Sie Ihre Algorithmen gegenseitig aus: eine Person liest vor, die andere führt die Aktionen aus, ohne nachzufragen.

\end{enumerate}



% ============================================================
\section*{Eigenschaften eines Algorithmus}

\begin{tcolorbox}[colback=white, colframe=black, boxrule=0.3pt]
\textbf{Definition:}\par
\vspace{2cm}

\textbf{Eigenschaften:}
\vspace{4.5cm}
\end{tcolorbox}

% ============================================================
\newpage
\section*{Quersumme}

Die \textbf{Quersumme} einer Zahl ist die Summe ihrer Ziffern. Beispiele: Die Quersumme von 44 ist $4+4 = 8$. Die Quersumme von 2026 ist $2+0+2+6 = 10$. 

\begin{enumerate}
\item Beschreiben Sie einen Algorithmus, der die Quersumme einer natürlichen Zahl $n$ berechnet. 

\begin{tcolorbox}[colback=white, colframe=black, boxrule=0.3pt, height=7cm]
\vspace*{\fill}
\end{tcolorbox}

\item Implementieren Sie anschließend Ihren Algorithmus als eine Python-Funktion (siehe Aufgabe im ISIS-Kurs).
\end{enumerate}
\end{document}
